\section{A Full Day's Wage for One Hour's Work}
\label{sec:opening}

At the centre of this Sunday's Mass stands a parable that ends in a complaint
and a question. A householder goes out early in the morning to hire labourers into
his vineyard and agrees with them ``for a penny a day'', the \latin{denarius}
of the Latin text. He goes out again about the third hour, the sixth and the
ninth, promising ``what shall be just'', and once more about the eleventh. When
evening comes the steward pays the last first, and every man receives a
denarius. Those who came at dawn murmur: ``These last have worked but one
hour, and thou hast made them equal to us, that have borne the burden of the
day and the heats.'' The householder answers one of them: ``Friend, I do thee
no wrong: didst thou not agree with me for a penny? Take what is thine, and go
thy way: I will also give to this last even as to thee. Or, is it not lawful
for me to do what I will? Is thy eye evil, because I am good?'' The Lord
concludes: ``So shall the last be first and the first last.''

{\small\noindent Scripture is quoted throughout in the Douay--Rheims version as
revised by Bishop Challoner, a public-domain study translation made from the
Latin Vulgate. It is not the text proclaimed at Mass: the English heard in the
celebration is that of the \work{Lectionary for Mass for Use in the Dioceses of
the United States of America}, which alone is the approved wording. Psalms are
numbered as the Lectionary numbers them, with the Vulgate number, which the
Douay--Rheims and the Latin books use, in parentheses.\par}
\medskip

Around the parable stand texts of other books and other ages: the antiphons
and prayers that the \work{Roman Missal}, Third Edition, gives under the
Twenty-fifth Sunday in Ordinary Time, \latin{Dominica XXV ``per annum''}, and
the readings of Year~A, no.~133 of the \work{Lectionary for Mass}. In the first
reading Isaiah cries, ``Seek ye the Lord, while he may be found'', and
promises the sinner who returns a God ``bountiful to forgive''; the reason he
gives is the sentence ``my thoughts are not your thoughts: nor your ways my
ways''. The psalm answers that the Lord is ``gracious and
merciful'', ``just in all his ways'' and ``nigh unto all them that call upon
him''. In the second reading Paul, a prisoner, weighs death, which would
bring him to Christ, against life, which means ``the fruit of labour'' for
his churches, and judges that to remain ``is needful for you''. Before the
Gospel the assembly asks the Lord to open its heart, in a verse adapted from
the story of Lydia, the first convert at Philippi. The Missal's own texts,
which serve this Sunday in all three years of the cycle, frame the readings
with a promise and three petitions: an oracle in which the Lord names himself
the salvation of his people, a Collect that rests the whole law on the love of
God and of neighbour, a prayer that what faith professes the sacrament may
give, and a prayer that the effect of redemption be held in the mysteries and
in conduct. At Communion the Missal offers either the psalmist's wish that his
ways be directed to keep the commandments or the Good Shepherd's word that he
knows his own.

The \work{Ordo lectionum Missae} marks what it singles out in each reading by
the title it sets over it. Over the first reading stands \latin{Non sunt
cogitationes meae cogitationes vestrae} (Isa 55:8); over the Gospel, \latin{An
oculus tuus nequam est, quia ego bonus sum?} (Mt 20:15); over the second
reading, \latin{Mihi vivere Christus est} (Phil 1:21). Its introduction
explains that in Ordinary Time the Old Testament reading is chosen for its
relation to the Gospel and that the titles show the relation, while the
Apostle and the Gospel are each read in their own semi-continuous course
(\work{Praenotanda} 105--107). The Lectionary thus sets God's thoughts, which
are not ours, beside the householder's question about his own goodness. The
letter to the Philippians begins on this Sunday a course of four Sundays.

The Sunday's question is therefore what kind of justice pays the last as the
first, and what the complaint against it reveals. The Fathers and Doctors who
expounded these passages answer along three lines, and each line, carried
through every text of the Mass, yields an interpretation of the whole. The first takes
the householder's day as the history of salvation: God has called labourers
from the first just man to the last, Israel through the long day and the
nations at its end, and gives them all the one wage of eternal life. The
second takes the day as a single human life: God calls at every age, the
late comer is truly received, and because no one knows the length of his day
no one may put off the call. The third begins from the evening and the
householder's reply: what the last receive is a gift, the first are wronged
in nothing, and the wage is finally the giver himself. What its witnesses say
of their own passages supplies each: the parable, the prophecy,
the two psalms, the letter, the story of Lydia and the discourse of the
shepherd. None of these witnesses expounds this Mass as a whole. All the
appointed elements enter each interpretation, and its own literal,
allegorical, moral and anagogical senses close it. What the three hold in
common, and where their answers part, comes last, in a short comparison.

\subsection{The appointed elements at a glance}

\begin{maptable}
Entrance Antiphon & \latin{Salus populi ego sum}; a composed text whose basis
the official Antiphonary identifies as Ps 37 (36):39--40, though the Missal
prints no locator &
The Lord names himself the salvation of the people, promises to hear them
from whatever tribulation they cry out of, and declares himself their Lord
for ever.\\
Collect & \latin{Deus, qui sacrae legis omnia constituta} &
The address rests all that the sacred law lays down on the love of God and
of neighbour; the petition is that those who keep God's precepts may reach
eternal life.\\
First Reading & Isa 55:6--9 &
Seek the Lord while he may be found; let the wicked return, for God is
bountiful to forgive; his thoughts and ways are as far above ours as heaven
above earth.\\
Responsorial Psalm & Ps 145 (144):2--3, 8--9, 17--18; response from v.~18a &
Daily praise of the Lord's greatness; the Lord gracious, merciful and patient,
sweet to all; just in all his ways and near to all who call on him in truth.\\
Second Reading & Phil 1:20c--24, 27a &
Christ will be magnified in Paul's body by life or by death; to live is
Christ and to die is gain; he is torn between the two and stays for his
people's sake; let their life be worthy of the Gospel.\\
Alleluia verse & \latin{Aperi, Domine, cor nostrum}; cf.\ Acts 16:14b, an
adaptation &
A petition that the Lord open our heart to heed the words of his Son, formed
from the account of Lydia's conversion.\\
Gospel & Mt 20:1--16a &
The labourers hired at five hours of the day and paid alike; the murmuring of
the first; the householder's reply; the last first and the first last.\\
Prayer over the Offerings & \latin{Munera, quaesumus, Domine, tuae plebis} &
God is asked to accept his people's gifts, so that what they profess in faith
they may lay hold of through the heavenly sacraments.\\
Communion Antiphon, first option & \latin{Tu mandasti mandata tua custodiri
nimis}; Ps 119 (118):4--5 &
God has commanded his commandments to be kept most diligently; the speaker
wishes that his ways be directed to keep them.\\
Communion Antiphon, second option & \latin{Ego sum pastor bonus}; Jn 10:14 &
``I am the good shepherd: and I know mine, and mine know me.''\\
Prayer after Communion & \latin{Quos tuis, Domine, reficis sacramentis} &
Those whom God restores with his sacraments are commended to his continual
help, that they may possess the effect of redemption both in the mysteries
and in their conduct.\\
\end{maptable}

\noindent The two Communion antiphons are alternatives: the Missal prints the
second under \latin{Vel}, and one of them is used. Every other element is
required. The formulary has no Sequence, Offertory antiphon, proper Preface,
prayer over the people or solemn blessing, and the Lectionary gives no shorter
form of the Gospel.
